\documentclass[../DoAn.tex]{subfiles}
\begin{document}

\begin{center}
    \Large{\textbf{ABSTRACT}}\\
\end{center}
\vspace{1cm}

Music discovery on current platforms relies mostly on metadata (artist, genre) or on user listening history, which only works well when listeners already know what they want. Listeners, however, often look for music by emotion, yet ordinary users find it hard to put that emotion into precise words. Existing approaches either depend on user-behaviour data and thus face the cold-start problem, or rely on hand-crafted audio features that capture emotion weakly; existing colour-to-music systems, in turn, are tied to Western data and map colours with manual rules. For Vietnamese music the difficulty is greater still, as no human-labelled Vietnamese music-emotion dataset exists for direct training.

Starting from the observation that picking a colour is far easier than naming an emotion, this thesis builds Brightify, whose primary feature is colour-based music recommendation, complemented by a similar-song feature. The chosen approach maps both colours and songs into a single two-dimensional Valence--Arousal space (Russell's circumplex) --- the theoretical bridge that lets colour and music meet, since colour--music links are mediated by emotion. Each song is assigned an offline-computed emotion coordinate: valence comes mainly from lyrics through an ensemble of grounded textual signals (the Vietnamese NRC-VAD lexicon, ViSoBERT, and XLM-R on EmoBank), while arousal comes from the self-supervised MuQ audio model combined with tempo and loudness; colours are mapped into the same space via Oklab, the ICEAS norms, and Whiteford's colour--music relationship. The system retrieves songs nearest in emotion while enforcing acoustic coherence among the results.

The main contributions are a scientifically grounded, reproducible emotion-labelling pipeline for Vietnamese music (no fine-tuning, and serving uses only precomputed-feature lookups, so it is deterministic); an emotion-based colour-to-music mechanism over a shared learned space; and a rigorous multi-layer evaluation with a model bake-off. On a catalogue of $5{,}138$ songs, the system outperforms strong baselines on offline metrics with statistical testing (FDR correction, bootstrap confidence intervals) and independent checks (true tempo, PMEmo cross-domain transfer, an independent GPT emotion reference). As no real user study has been conducted yet, these results are offline metrics.

\end{document}
